% GENERATED FILE. Edit canonical lessons, then run npm run generate:books.
% canonical-source-hash: fnv1a64:e71662db16667896
% canonical-lessons: LA-C508-ignavia, LA-C508-pigritia, LA-C508-segnis, LA-C508-strenuus, LA-C508-assiduus

\chapter{Laziness and energy}
\label{ch:la-laziness-and-energy}
\hlchaptermodality{508}

\begin{chapteropening}
\textbf{By the end of this chapter:} \emph{I can describe lazy and hard-working people.}

Complete the last lesson of chapter 508: I can describe lazy and hard-working people.
\end{chapteropening}

\section[ignāvia, ignāviae]{ignāvia, ignāviae — cowardice; laziness}

\label{lesson:LA-C508-ignavia}



\begin{quote}
Before the new one: say the Latin for insolent, then the Latin for boastfulness.
\end{quote}

\subsection*{You'll want to know: ignāvia, ignāviae}
\textbf{ignāvia, ignāviae} — \textquotedblleft{}cowardice; laziness\textquotedblright{}.

\begin{quote}
Ignāvia mīlitum ducem īrātum facit.
\end{quote}

\begin{quote}
\emph{The soldiers' cowardice makes the general angry.}
\end{quote}

Feminine, first declension, from \emph{ignāvus} — lazy, cowardly.

\begin{cousinweb}
\textbf{ignāvia} is built on \emph{ignāvus} \textquotedblleft{}idle, cowardly\textquotedblright{}, from \emph{in-} \textquotedblleft{}not\textquotedblright{} and old \emph{gnāvus} \textquotedblleft{}busy, active\textquotedblright{} (Lewis \& Short).
\end{cousinweb}

\begin{grammarlens}[title={What you've built}]
One more word for telling, explaining and describing at B1.
\end{grammarlens}

\subsection*{Your turn}
\begin{itemize}
\raggedright
  \item \emph{Say it:} \emph{ignāvia, ignāviae}
  \item \emph{Say it:} \emph{ignāvia, ignāviae}, once more
  \item \emph{Say it:} a sentence of your own with \emph{ignāvia, ignāviae}
  \item \emph{Recall:} say the Latin for insolent, then the Latin for boastfulness, then say \emph{ignāvia, ignāviae} again
\end{itemize}

\begin{tcolorbox}[breakable,colback=teal!4,colframe=teal!35!black,title={Before you move on}]
What does ignāvia, ignāviae mean? (\textquotedblleft{}Cowardice; laziness\textquotedblright{}.) Say it once more.
\end{tcolorbox}
\section[pigritia, pigritiae]{pigritia, pigritiae — sloth, laziness}

\label{lesson:LA-C508-pigritia}



\begin{quote}
Before the new one: say the Latin for boastfulness, then the Latin for cowardice.
\end{quote}

\subsection*{You'll want to know: pigritia, pigritiae}
\textbf{pigritia, pigritiae} — \textquotedblleft{}sloth, laziness\textquotedblright{}.

\begin{quote}
Propter pigritiam servus nihil facit.
\end{quote}

\begin{quote}
\emph{Because of his laziness the slave does nothing.}
\end{quote}

Feminine, first declension, from \emph{piger} — lazy.

\begin{cousinweb}
\textbf{pigritia} is built on \emph{piger} \textquotedblleft{}lazy, slow\textquotedblright{} (Lewis \& Short; EtymDB), which EtymDB links with \emph{pigeō} \textquotedblleft{}to feel annoyance\textquotedblright{}.
\end{cousinweb}

\begin{grammarlens}[title={What you've built}]
One more word for telling, explaining and describing at B1.
\end{grammarlens}

\subsection*{Your turn}
\begin{itemize}
\raggedright
  \item \emph{Say it:} \emph{pigritia, pigritiae}
  \item \emph{Say it:} \emph{pigritia, pigritiae}, once more
  \item \emph{Say it:} a sentence of your own with \emph{pigritia, pigritiae}
  \item \emph{Recall:} say the Latin for boastfulness, then the Latin for cowardice, then say \emph{pigritia, pigritiae} again
\end{itemize}

\begin{tcolorbox}[breakable,colback=teal!4,colframe=teal!35!black,title={Before you move on}]
What does pigritia, pigritiae mean? (\textquotedblleft{}Sloth, laziness\textquotedblright{}.) Say it once more.
\end{tcolorbox}
\section[segnis, segne]{segnis, segne — sluggish, slow}

\label{lesson:LA-C508-segnis}



\begin{quote}
Before the new one: say the Latin for cowardice, then the Latin for sloth.
\end{quote}

\subsection*{You'll want to know: segnis, segne}
\textbf{segnis, segne} — \textquotedblleft{}sluggish, slow\textquotedblright{}.

\begin{quote}
Asinus segnis lentē ambulat.
\end{quote}

\begin{quote}
\emph{The sluggish donkey walks slowly.}
\end{quote}

Two endings. \emph{Nōn segniter} — not slackly, with energy.

\begin{cousinweb}
\textbf{segnis} is a native Latin adjective; Lewis \& Short reports that it is commonly linked with \emph{sequor} \textquotedblleft{}to follow\textquotedblright{}, but the origin is not settled.
\end{cousinweb}

\begin{grammarlens}[title={What you've built}]
One more word for telling, explaining and describing at B1.
\end{grammarlens}

\subsection*{Your turn}
\begin{itemize}
\raggedright
  \item \emph{Say it:} \emph{segnis, segne}
  \item \emph{Say it:} \emph{segnis, segne}, once more
  \item \emph{Say it:} a sentence of your own with \emph{segnis, segne}
  \item \emph{Recall:} say the Latin for cowardice, then the Latin for sloth, then say \emph{segnis, segne} again
\end{itemize}

\begin{tcolorbox}[breakable,colback=teal!4,colframe=teal!35!black,title={Before you move on}]
What does segnis, segne mean? (\textquotedblleft{}Sluggish, slow\textquotedblright{}.) Say it once more.
\end{tcolorbox}
\section[strēnuus, strēnua, strēnuum]{strēnuus, strēnua, strēnuum — vigorous, active}

\label{lesson:LA-C508-strenuus}



\begin{quote}
Before the new one: say the Latin for sloth, then the Latin for sluggish.
\end{quote}

\subsection*{You'll want to know: strēnuus, strēnua, strēnuum}
\textbf{strēnuus, strēnua, strēnuum} — \textquotedblleft{}vigorous, active\textquotedblright{}.

\begin{quote}
Agricola strēnuus ā prīmā lūce labōrat.
\end{quote}

\begin{quote}
\emph{The vigorous farmer works from first light.}
\end{quote}

Brisk and energetic, in work or in war. English \emph{strenuous}.

\begin{cousinweb}
\textbf{strēnuus} is a native Latin adjective; Lewis \& Short and EtymDB compare Greek \emph{stereos} \textquotedblleft{}firm, hard\textquotedblright{}, but the link is not settled. English \emph{strenuous} comes from it.
\end{cousinweb}

\begin{grammarlens}[title={What you've built}]
One more word for telling, explaining and describing at B1.
\end{grammarlens}

\subsection*{Your turn}
\begin{itemize}
\raggedright
  \item \emph{Say it:} \emph{strēnuus, strēnua, strēnuum}
  \item \emph{Say it:} \emph{strēnuus, strēnua, strēnuum}, once more
  \item \emph{Say it:} a sentence of your own with \emph{strēnuus, strēnua, strēnuum}
  \item \emph{Recall:} say the Latin for sloth, then the Latin for sluggish, then say \emph{strēnuus, strēnua, strēnuum} again
\end{itemize}

\begin{tcolorbox}[breakable,colback=teal!4,colframe=teal!35!black,title={Before you move on}]
What does strēnuus, strēnua, strēnuum mean? (\textquotedblleft{}Vigorous, active\textquotedblright{}.) Say it once more.
\end{tcolorbox}
\section[assiduus, assidua, assiduum]{assiduus, assidua, assiduum — diligent, unremitting}

\label{lesson:LA-C508-assiduus}



\begin{quote}
Before the new one: say the Latin for sluggish, then the Latin for vigorous.
\end{quote}

\subsection*{You'll want to know: assiduus, assidua, assiduum}
\textbf{assiduus, assidua, assiduum} — \textquotedblleft{}diligent, unremitting\textquotedblright{}.

\begin{quote}
Discipulus assiduus semper in scholā est.
\end{quote}

\begin{quote}
\emph{The diligent pupil is always at school.}
\end{quote}

From \emph{assideō} — to sit beside: someone who stays at the task. English \emph{assiduous}.

\begin{cousinweb}
\textbf{assiduus} is built on \emph{assideō} \textquotedblleft{}to sit beside\textquotedblright{}, from \emph{ad-} and \emph{sedeō} \textquotedblleft{}to sit\textquotedblright{} (Lewis \& Short; EtymDB). English \emph{assiduous} comes from it.
\end{cousinweb}

\begin{grammarlens}[title={What you've built}]
One more word for telling, explaining and describing at B1.
\end{grammarlens}

\subsection*{Your turn}
\begin{itemize}
\raggedright
  \item \emph{Say it:} \emph{assiduus, assidua, assiduum}
  \item \emph{Say it:} \emph{assiduus, assidua, assiduum}, once more
  \item \emph{Say it:} a sentence of your own with \emph{assiduus, assidua, assiduum}
  \item \emph{Recall:} say the Latin for sluggish, then the Latin for vigorous, then say \emph{assiduus, assidua, assiduum} again
\end{itemize}

\begin{tcolorbox}[breakable,colback=teal!4,colframe=teal!35!black,title={Before you move on}]
What does assiduus, assidua, assiduum mean? (\textquotedblleft{}Diligent, unremitting\textquotedblright{}.) Say it once more.
\end{tcolorbox}
